\documentclass[11pt]{article}
\usepackage{amsmath,amssymb,amsthm}
\usepackage[margin=1.2in]{geometry}
\usepackage{hyperref}

\newtheorem{theorem}{Theorem}
\newtheorem{lemma}[theorem]{Lemma}
\newtheorem{proposition}[theorem]{Proposition}
\newtheorem{corollary}[theorem]{Corollary}
\theoremstyle{definition}
\newtheorem{definition}[theorem]{Definition}
\newtheorem{remark}[theorem]{Remark}
\newtheorem{question}[theorem]{Question}
\newtheorem{conjecture}[theorem]{Conjecture}

\newcommand{\Z}{\mathbb{Z}}
\newcommand{\N}{\mathbb{N}}
\newcommand{\SA}{S_{\mathrm{A}}}
\newcommand{\pos}{\mathrm{pos}}
\newcommand{\OG}{\mathrm{OG}}

\title{Structural rigidity in the Erd\H{o}s--Graham two-set permutation problem}
\author{Will Kasel\thanks{With extensive machine assistance from Claude (Anthropic);
SAT certificates and code at [repository]. All proofs have been verified by hand
unless explicitly marked as machine-checked.}}
\date{Draft: August 2026}

\begin{document}
\maketitle

\begin{abstract}
A sequence $a_1, a_2, \dots$ \emph{contains a monotone 3-AP} if some three terms in
increasing or decreasing positions form an arithmetic progression in the matching
order. Davis, Entringer, Graham and Simmons (1977) proved that every permutation of
$\Z^+$ contains a monotone 3-AP, and Erd\H{o}s and Graham (1979) asked whether
$\Z^+$ can be partitioned into two sets, each of which admits a permutation with no
monotone 3-AP (three sets suffice; problem \#197 of the Erd\H{o}s problem
collection). We prove several structural results constraining any solution.
(i) A permutable set admits no infinite doubling orbit $u_{k+1}=2u_k-f_k$ with
$f_k$ drawn from a fixed finite subset. (ii) In any 3-AP-free arrangement, every
newly placed value must be balanced: the number of already-placed values in
$[1,v)$ and in $(v,2v)$ agree up to the codensity of the set above $v$.
(iii) No 3-AP-free arrangement of $\{k,2k,3k,5k\}\subseteq S$ places all
ratio-$\ge 2$ pairs in increasing order. (iv) For the canonical dyadic partition,
no solution can play any cofinite family of dyadic blocks as contiguous runs: a
halving self-reduction plus finitely many machine-checked base cases shows the
associated ``zone'' systems are infeasible, and a finish-time descent completes
the argument. (v) In contrast, all natural finite fragments of the problem are
satisfiable, and we give an exact characterization of solvability in terms of
bounded-delay towers of finite arrangements, which pinpoints the problem's
content in the growth of a single computable function.
(vi) For the dyadic partition we prove an unconditional reduction to a single
scale-uniform statement: an ``order gadget'' $\OG(M)$ on the interval
$(M, 2M]$ --- no monotone in-block $3$-AP, plus fifteen guard-precedence
axioms induced by the values $15$ and $16$ --- whose infeasibility at
infinitely many dyadic scales $M = 2^{2t-1}$ implies that $\SA$ admits no
valid permutation.
(vii) We then prove, by hand, that infeasibility: for every
$M \equiv 0 \pmod 8$ with $M \ge 16$, already the three-axiom core
$\mathrm{C3} \subset \OG(M)$ is inconsistent with AP-freeness. The proof
runs on a small toolkit of ladder lemmas (zigzag propagation, phase
dichotomy, a transfer lock, and a mirror-flood induction) and exhibits the
mod-$8$ arithmetic as a residue condition on two flood centers; every step
is machine-verified at $100$ scales up to $M = 1024$ and independently
cross-validated. Combining (vi) and (vii):
\textbf{$\SA$ admits no permutation free of monotone $3$-APs} --- the
canonical dyadic partition does not resolve the Erd\H{o}s--Graham problem.
We also prove a clean
absorption theorem for van der Corput orderings of odd residue classes which
explains the fractal structure observed in all satisfying fragments.
\end{abstract}

\section{Introduction}

% TODO: history: DEGS77, ErGr79/80 (p. 338: "very annoying question"),
% LeSaulnier--Vijay 2011 (alpha/beta bounds + conjecture), Geneson 2018,
% erdosproblems.com #197.

\section{Definitions and conventions}

\begin{definition}
A set $S \subseteq \Z^+$ is \emph{permutable} if there is a bijection
$\pi\colon \N \to S$ such that no $i<j<k$ have $\pi(i),\pi(j),\pi(k)$ forming an
arithmetic progression with $\pi(i)<\pi(j)<\pi(k)$ or $\pi(i)>\pi(j)>\pi(k)$.
\end{definition}

The completion of an ordered pair: if $x$ is placed before $y$, the value
$z = 2y - x$ (which continues the progression through $x,y$ in the direction of
$y$) must not be placed after $y$; equivalently every such $z \in S$ placed later
yields a monotone 3-AP. We use throughout the \emph{dyadic blocks}
$B_k = (2^{k-1}, 2^k]$ and the canonical two-set candidate
$\SA = \bigcup_{k \text{ even}} B_k$, $S_B = \Z^+ \setminus \SA$.

\section{The orbit obstruction}

\begin{lemma}[Orbit obstruction]\label{lem:orbit}
Let $S$ be permutable and $F \subseteq S$ finite. There is no infinite sequence
$u_0 < u_1 < \cdots$ in $S$ with $u_{k+1} = 2u_k - f_k$, $f_k \in F$.
\end{lemma}

\begin{proof}
Let $q$ be the largest position of an element of $F$ in the permutation. Since
the $u_k$ are distinct, at most $q$ of them occupy positions $\le q$; beyond the
largest such index, each pair $(f_k, u_k)$ is increasingly placed, so its
completion $u_{k+1} \in S$ must be placed before $u_k$, giving an infinite strictly
decreasing sequence of positions.
\end{proof}

Taking $S = \Z^+$, $F = \{1,2\}$ recovers the theorem of Davis--Entringer--Graham--Simmons.

\section{The balance law}

\begin{lemma}[Balanced placement]\label{lem:balance}
Let $S$ be permutable via $\pi$, and consider the moment a value $v$ is placed.
Let $L$ (resp.\ $H$) be the number of already-placed values in $[1,v)$ (resp.\
$(v,2v)$). Then $|H - L| \le |S^c \cap (v, 2v)| + O(1)$.
\end{lemma}

\begin{proof}
Each placed $x < v$ forms an increasingly placed pair $(x,v)$ with completion
$2v - x \in (v, 2v)$, distinct for distinct $x$; each completion lying in $S$
must already be placed. This gives $H \ge L - |S^c \cap (v,2v)|$. Symmetrically,
each placed $u \in (v, 2v)$ forms a decreasingly placed pair $(u,v)$ with
completion $2v - u \in (0, v)$, whence $L \ge H - |S^c \cap (v,2v)| - O(1)$.
\end{proof}

\begin{corollary}[Records]\label{cor:records}
When a left-to-right maximum $v$ is placed, at most $|S^c \cap (v,2v)| + O(1)$
smaller values are already placed.
\end{corollary}

\section{Two absorption lemmas}

\begin{lemma}[Ratio-ascent obstruction]\label{lem:R}
Let $S \supseteq \{k, 2k, 3k, 5k\}$. No 3-AP-free arrangement of $S$ places $w$
before $u$ for every pair $u \ge 2w$ (elements of $S$).
\end{lemma}

\begin{proof}
The hypothesis forces $k \prec 2k$, $k \prec 3k$, $2k \prec 5k$. The progression
$(k,2k,3k)$ with $k \prec 2k$, $k \prec 3k$ forces $3k \prec 2k$; then
$3k \prec 2k \prec 5k$, and $(k,3k,5k)$ is an increasing 3-AP.
\end{proof}

\begin{theorem}[van der Corput absorption]\label{thm:vdc}
Order the odd residues $c$ modulo $2^k$ by the bit-reversal (van der Corput) rank
of $(c-1)/2 \bmod 2^{k-1}$. Then for any two classes $a$ ordered before $b$, the
class $2b - a \bmod 2^k$ is ordered strictly before $b$.
\end{theorem}

\begin{proof}
Write $\alpha = (a-1)/2$, $\beta = (b-1)/2$; the completion class has parameter
$\zeta = 2\beta - \alpha \bmod 2^{k-1}$. If the lowest bits of $\alpha,\beta$
agree, then so does $\zeta$'s, and dividing by 2 preserves the form
$\zeta' = 2\beta' - \alpha'$. At the first bit where they differ, the assumption
$\mathrm{rev}(\alpha) < \mathrm{rev}(\beta)$ forces $\alpha$'s bit to be $0$ and
$\beta$'s to be $1$; since $\zeta \equiv \alpha$ at that bit, $\zeta$'s bit is
$0$, so $\mathrm{rev}(\zeta) < \mathrm{rev}(\beta)$ strictly.
\end{proof}

\section{Contiguous blocks cannot work}

\begin{definition}
The \emph{zone system} $Z(M)$: arrange the interval $(M, 2M]$ so that (a) no
monotone 3-AP occurs, and (b) for every $y < z$ in $(M,2M]$ with
$2y - z \in (M/4, M/2]$, $z$ precedes $y$.
\end{definition}

Condition (b) is forced when the block $(M/4, M/2]$ is entirely placed before a
contiguous run of $(M, 2M]$ within a team's sequence (its elements' completions
$2y - x$ then must precede $y$).

\begin{lemma}[Halving descent]\label{lem:halving}
If $Z(M)$ is satisfiable then so is $Z(\lfloor M/2 \rfloor)$.
\end{lemma}

\begin{proof}
Restrict a witness to the even values of $(M,2M]$ and halve; arithmetic
progressions and the zone constraints scale exactly (the boundary cases are
routine), and a superfluous top element may be deleted.
\end{proof}

\begin{lemma}[Base cases; machine-checked]\label{lem:bases}
$Z(M)$ is unsatisfiable for $16 \le M \le 31$.
\end{lemma}

\begin{theorem}[No contiguous-run solutions]\label{thm:contig}
In any 3-AP-free permutation of $\SA$, infinitely many blocks $B_{2k}$ are
\emph{not} placed as contiguous runs.
\end{theorem}

\begin{proof}
Suppose cofinitely many blocks are contiguous runs. Machine-checked
computations (with Lemma~\ref{lem:halving} extending them to all $M \ge 16$;
see also the run-tolerance bounds in Section~7) show that when block $B_{2k}$ is
played as a run, all but $O(1)$ elements of $B_{2k-2}$ must be unplaced at that
time, so $B_{2k-2}$ finishes after $B_{2k}$; iterating yields an infinite
strictly decreasing sequence of finish positions.
\end{proof}

\section{Local realizability and the tower characterization}

\begin{definition}[Pure-complete systems]
For an $\SA$-block top $X$, the \emph{pure-complete-$X$ system} asks for an
arrangement of $\SA \cap [1, X]$ with no monotone 3-AP. (All completions of
in-range pairs lie in $(X, 2X]$, an odd block, or inside the range; hence this
is exactly the restriction of the infinite problem.)
\end{definition}

\begin{lemma}[Horizon collapse]
Deleting any set of values from a valid arrangement preserves validity; hence
auxiliary values above $X$ never aid the pure-complete-$X$ system.
\end{lemma}

\begin{theorem}[Tower characterization]
$\SA$ is permutable iff there exists $B\colon \N \to \N$ such that for every
block top $X$ the system [pure-complete-$X$ and every $v$ has at most $B(v)$
predecessors] is satisfiable.
\end{theorem}

\begin{proof}
($\Rightarrow$) restrict a permutation and take $B(v)$ its positions.
($\Leftarrow$) the restriction sets are finite and nested; K\H{o}nig's lemma
yields a consistent tower whose limit has each $v$ preceded by at most $B(v)$
values, hence order type $\omega$; all constraints are finitary and inherited.
\end{proof}

Computationally: pure-complete-$X$ is satisfiable for $X \le 1024$ (machine
certificates), and the completion-pressure function
$g_X(L) = \min \max_{v \le L} \pos(v)$ satisfies $g_{256}(64) = 153$
(optimal): eighty-seven percent of the block $(128,256]$ must precede the
completion of $[1,64]$, with the residue concentrated in a single residue
class modulo 8 --- the reservoir structure that any construction must
implement. Divergence of $g_X(L)$ in $X$ for fixed $L$ would prove
non-permutability; the mechanism is blocked by the measured subset-tolerance
(near-total), and $g$ is expected to stabilize.

\subsection{Self-similar witnesses, fragility, and pumping obstructions}

Adding the exact scale-invariance $u \prec w \iff 4u \prec 4w$ to the
pure-complete systems remains satisfiable at $X = 256$ and $X = 1024$
(independently audited witnesses); however the audited $1024$-witness does not
extend to a self-similar $4096$-arrangement (machine-certified infeasibility),
so the self-similar tower also has dead branches.

We also record a calibration fact: requiring, in addition, that the $\pm 1$
neighbors (within the same block and team) of each forced completion also
precede the pair's later element (``radius-1 robustness'') is infeasible in
every setting we tested --- including plain intervals $[1, n]$ from $n = 8$.
Radius-1 robustness is thus unachievable generically in three-AP arrangement
problems and cannot serve as a satisfiability criterion; in particular it does
not discriminate between partitions, and pumping arguments must instead manage
rounding drift on the thin family of far-pair constraints only. What does
follow is the modest but genuine statement that no placement rule in which
position varies slowly with value (Lipschitz-type rules) can produce valid
arrangements beyond bounded scales.

\section{Partition rigidity}

\begin{lemma}[Ray piercing]\label{lem:ray}
Let $\Z^+ = A \sqcup B$ with both parts permutable. Then for every $a \in A$ and
every $d \ge 1$, the set $\{k \ge 0 : a + 2^k d \in B\}$ is infinite, and
symmetrically with $A$ and $B$ exchanged.
\end{lemma}

\begin{proof}
If $a + 2^k d \in A$ for all $k \ge k_0$, the values $u_j = a + 2^{k_0+j} d$
form an infinite orbit $u_{j+1} = 2u_j - a$ inside $A$ with $F = \{a\}$,
contradicting Lemma~\ref{lem:orbit}.
\end{proof}

Both teams must therefore alternate along every affine doubling ray; partitions
such as ``even 2-adic valuation vs.\ odd'' die immediately, and viable partitions
are forced toward dyadic-interval alternation, motivating the canonical
candidate $\SA$.

\section{Toward the main theorem: suffix-stacked class deferral}

% DRAFT SECTION — the construction program in its corrected form.

\begin{definition}[Deferral states]
For a scale $M$ (a power of two) let $D_M \subseteq (M/2, M]$ denote a single
residue class modulo $m(M)$ (class coherence; the witnesses select
$\{v \equiv 2 \bmod 2^{k/2}\}$). A \emph{deferral state} at scale $M$ is a
doom-free arrangement of $\SA \cap [1, 2M]$ of the suffix-stacked form
$[\,\mathrm{bulk} \mid D_{M/2} \mid D_{M}\,]$: the bulk contains everything
except the two most recent deferred classes, which are appended in scale
order.
\end{definition}

Machine-verified facts (certificates in the repository): deferral states exist
for $M = 32, 64, 128$; the suffix must be stacked in scale order (the reversed
stacking contains a Lemma~\ref{lem:R} configuration); no-extras variants are
infeasible at every scale (lookahead into the next block is necessary); any
single residue class serves as the deferred suffix, while generic sets of the
same size do not.

\begin{proposition}[Class-sink algebra]
Fix a modulus $m$ and a target class $r$. The attacker map
$\varphi_r(b) = 2b - r$ on $\Z/m\Z$ generates a tree rooted at $r$; hence in a
class-phase arrangement a single class can be scheduled last (all its
attacker-demands point inward), whereas a fully placed zone imposes root
demands at every class simultaneously, which close into cycles. For odd $r$
the shifted map has genuine cycles, so an odd class cannot be scheduled last
within the pure class calculus; deferred classes must be even, and the
$2\cdot\mathrm{odd}$ phase requires range truncation.
\end{proposition}

\section{The chunk reduction}\label{sec:chunk}

Every permutation of $\SA$ of order type $\omega$ decomposes into consecutive
finite segments (\emph{chunks}); conversely, a choice of finite fibers and
internal orders assembles into such a permutation. The following exact
reformulation is the coordinate system for everything below.

\begin{definition}
A \emph{stage function} is any $s\colon \SA \to \N$ with finite fibers.
Given $s$ and a linear order on each fiber, let $T$ be the concatenation.
\end{definition}

\begin{theorem}[Chunk reduction; exact]\label{thm:chunk}
$T$ avoids monotone $3$-term APs if and only if:
\begin{enumerate}
\item[(A)] no AP triple $(x,y,z)$, $z=2y-x$, in $\SA$ has strictly monotone
stages ($s(x)<s(y)<s(z)$ or $s(z)<s(y)<s(x)$); and
\item[(B)] each fiber order contains the \emph{forced pairs}
\[
\begin{array}{ll}
s(x)<s(y)=s(z):\ z \prec y, &\quad s(z)<s(y)=s(x):\ x \prec y,\\
s(x)=s(y)<s(z):\ y \prec x, &\quad s(y)=s(z)<s(x):\ y \prec z,
\end{array}
\]
and is non-monotone on every within-fiber triple. (Triples whose endpoint
stages lie on one side of the middle's are automatically safe.)
\end{enumerate}
Hence $\SA$ is $3$-permutable if and only if some stage function satisfies
(A) and (B). The fiber problems are finite and mutually independent given
$s$.
\end{theorem}

\begin{lemma}[Normalization: running-max chunking]\label{lem:normal}
Every $3$-AP-free permutation $\pi$ of $\SA$ (of order type $\omega$)
admits a valid scheme --- a stage function with finite fibers satisfying
(A)$\wedge$(B), whose concatenation is $\pi$ itself --- with
$s(v) \ge \mathrm{block}(v)/2$ for every $v$.
\end{lemma}

\begin{proof}
Set $s(v) = \max_{q \le \pos(v)} \mathrm{block}(\pi(q))/2$ (the running
maximum of half-block-indices). Along positions $s$ is non-decreasing, so
its fibers are consecutive segments of $\pi$ and the concatenation (with
the fiber orders induced by $\pi$) is $\pi$; Theorem~\ref{thm:chunk} then
gives (A)$\wedge$(B). Since $\mathrm{block}(v)/2$ enters the maximum at
$v$'s own position, $s(v) \ge \mathrm{block}(v)/2$. Each fiber is finite:
$\pi$ is onto $\SA$, so a value of block index $> 2\sigma$ occurs at some
finite position, after which the running maximum exceeds $\sigma$.
\end{proof}

\noindent The point of Lemma~\ref{lem:normal} is that non-negative
\emph{displacement} (next section) is a harmless normalization, not an
automatic property: an arbitrary chunk decomposition can place a
high-block value in an early chunk. Restrictions of normalized schemes
stay normalized, which is what the divergence criterion below consumes.

\begin{theorem}[DEGS via chunks]\label{thm:degs}
$\Z^+$ itself is not $3$-permutable. \emph{Proof:} suppose $s$ satisfies
(A) with finite fibers and put $\sigma = s(1)$. For every $y$ with
$s(y) > \sigma$, the triple $(1, y, 2y-1)$ forces $s(2y-1) \le s(y)$ by (A).
The orbit $y_0,\, y_{i+1} = 2y_i - 1$ is strictly increasing; only finitely
many of its members can have stage $\le \sigma$ (those fibers are finite),
after which its stages are non-increasing, hence eventually confined to one
finite fiber — impossible for an infinite set. \qed
\end{theorem}

\noindent This recovers the Davis–Entringer–Graham–Simmons theorem in the
chunk coordinates and isolates what a permutable team must do: break every
such reflected orbit out of the team. The dyadic team $\SA$ breaks
single-reflector orbits in two steps (completions land in odd blocks); the
question is whether the multi-scale accounting below can be survived.

\begin{lemma}[Same block]\label{lem:sameblock}
In every AP triple of $\SA$ the middle $y$ and top $z$ lie in the same even
block. ($z<2y\le 2^{k+1}$ and block $k{+}1$ is odd.)
\end{lemma}

\begin{lemma}[Escape]\label{lem:escape}
If $y \ge 3\cdot 4^{s-1}$ (top quarter of block $2s$) and $x \le 4^{s-1}$,
then $2y-x > 4^{s}$: completions of (lower value, top-quarter value) pairs
land in the odd block above — out of team.
\end{lemma}

\begin{lemma}[Bottom-half characterization]\label{lem:bottomhalf}
Deferring $D_k \subseteq$ block $k$ by uniform delay is (A)-clean if
$D_k \subseteq (2^{k-1}, 3\cdot 2^{k-2}]$: the witness midpoints
$y=(x+z)/2$ of a bottom-half $z$ against lower-block $x$ satisfy
$2^{k-2} < y \le 2^{k-1}$, i.e.\ lie in the odd block below. Conversely a
deferred top-half $z$ has genuine in-team witnesses, forcing its witness
midpoints (bottom-quarter values) to defer with it.
\end{lemma}

\section{The negative atlas}\label{sec:negative}

\begin{theorem}[Block-granular death]\label{thm:blockgranular}
No stage function whose fibers are unions of whole even blocks satisfies
(A)$\wedge$(B). \emph{Proof:} the fatal core $F_b$ (forced pairs into block
$b$ from block $b{-}2$ below, plus block-$b$ triples) is UNSAT for $b=6$
(machine, instant), and $v \mapsto 4v$ embeds $F_b$ into $F_{b+2}$: the map
preserves $\SA$, sends block $j$ to block $j{+}2$, and preserves AP triples;
UNSAT lifts to supersystems. Any block-granular scheme has a fiber boundary
above block $4$ and therefore contains some $F_b$. \qed
\end{theorem}

\begin{theorem}[$\times 4$ restriction; window death is permanent]
\label{thm:restriction}
For the window class $\{s : s(v)-\mathrm{block}(v)/2 \in [0,w]\}$:
feasibility at horizon $4N$ implies feasibility at horizon $N$ (restrict
along $v \mapsto 4v$ and shift stages by one). Consequently the machine
verdicts below are permanent (hold at all larger horizons), and an infinite
solution in a window class would imply feasibility at every finite horizon.
\end{theorem}

\noindent Machine atlas (CP-SAT and CDCL cross-checks; certificates in the
repository): bottom-half relief of depth $1$ and depth $\le 2$ infeasible at
$N=256$ (and $1024$); \textbf{window-$2$ infeasible at $N=1024$} (CP-SAT, 746s; independently CDCL/unary-ladder, 1760s) — hence dead
at every horizon; window-$2$ at $N=256$ is feasible only via whole-block
merges (minimum total displacement $41$), whose skeleton is
$C_s = 2^{s}\cdot[3\cdot 2^{s-2}, 2^{s}]$, the top-quarter multiples of
$2^{s}$ — precisely the set singled out by Lemma~\ref{lem:escape}.

\medskip
The surviving candidate shapes for a YES are stage functions with unbounded
displacement (chunks mixing stragglers of unboundedly many blocks); the
surviving NO program is to close the gap between
Theorem~\ref{thm:blockgranular}/\ref{thm:restriction} and general finite-fiber
schemes via the per-block deferral calculus of
Lemma~\ref{lem:bottomhalf}. Both are finite-structure questions; the same
lemmas hold verbatim for $S_B$ (machine-verified mirror at $4096$).

\section{The displacement ladder}\label{sec:ladder}

For a stage function $s$ write $\delta(v) = s(v) - \mathrm{block}(v)/2$
(\emph{displacement}); a scheme is \emph{normalized} if $\delta \ge 0$
everywhere, which by Lemma~\ref{lem:normal} is no loss of generality for
schemes induced by genuine permutations. Define
\[
L(m) \;=\; \min_{\text{normalized schemes at horizon } 4^m}\ \max_{v \le 16} \delta(v),
\]
the minimum over all normalized finite-fiber schemes satisfying
(A)$\wedge$(B) at horizon $4^m$ (no bound on any other value's
displacement).

\begin{proposition}[Ladder rungs; machine-certified, solver-cross-checked]
$L(3) \le 1$; \ $L(4) = 2$ (infeasibility of the cap $\delta \le 1$ on
$\{v \le 16\}$ at $N=256$ verified independently by CP-SAT, 252s, and
CDCL with a unary-ladder encoding, 15s); \ $L(5) = 3$ (lower: CDCL, 2304s; upper: a window-3 witness at $N=1024$, 995s).
Moreover the minimal infeasible cap-coalition at $N=256$ is exactly
$\{15, 16\}$ (deletion-minimal): every scheme has $\delta(15) \ge 2$ or
$\delta(16) \ge 2$ — the crowns of the two competing class towers
($15 = 1111_2$, apex of the $\equiv -1$ skeleton; $16 = 10000_2$, the block
top, apex of the $\equiv 0$ skeleton).
\end{proposition}

\begin{theorem}[Divergence criterion]\label{thm:divergence}
If $L(m) \to \infty$ then $\SA$ is not $3$-permutable. \emph{Proof:} an
infinite valid permutation induces (Lemma~\ref{lem:normal}) a
\emph{normalized} valid scheme, and by restriction a normalized valid
scheme at every horizon $4^m$ (completions of values $\le 4^m$ that lie
beyond the horizon fall in the odd block $(4^m, 2\cdot 4^m]$, out of
team), and the ten displacements
$\delta(3), \delta(4), \dots, \delta(16)$ are constants of the infinite
scheme. For $m$ with $L(m) > \max_{v\le 16}\delta(v)$ this is a
contradiction, by the exactness of Theorem~\ref{thm:chunk}. \qed
\end{theorem}

The conjectured growth $L(m) = m-2$ matches the window-death schedule
(window $w$ dying at horizon $4^{w+3}$) and the observed per-block descent
of displacement in minimal solutions ($\delta$ profile $(m{-}2, m{-}3,
\dots, 1, 0)$ across blocks $4, 6, \dots, 2m$). The missing step is a
self-propagating squeeze $L(m{+}1) \ge L(m) + 1$; the $\times 4$
restriction alone yields only the diagonal transfer
$L_{b+1}(m{+}1) \ge L_b(m)$ for the capped set $[1, 4^b]$.

\section{The order gadget and the reduction}\label{sec:og}

The machinery of Sections~\ref{sec:chunk}--\ref{sec:ladder} compresses,
for the canonical candidate $\SA$, into a one-scale-at-a-time statement
about linear orders of single dyadic blocks. This section states that
reduction precisely and proves it unconditionally, and reports the machine
verification of its finite instances; Section~\ref{sec:c3} then proves the
required infeasibility by hand on the residue class $M \equiv 0 \pmod 8$
--- which contains every dyadic scale --- completing the proof of the main
theorem (Theorem~\ref{thm:main}). Throughout,
$b_j = M + j$ (\emph{bottoms}) and $t_i = 2M - i$ (\emph{tops}) denote
elements of the interval $(M, 2M]$.

\begin{definition}[Order gadget]\label{def:og}
For an integer $M \ge 16$, the \emph{order gadget} $\OG(M)$ asks for a
linear order $\prec$ of the interval $(M, 2M]$ such that
\begin{enumerate}
\item[(i)] \emph{(no monotone AP)} for every arithmetic progression
$a < b < c$ inside $(M, 2M]$, neither $a \prec b \prec c$ nor
$c \prec b \prec a$ holds; and
\item[(ii)] \emph{(guards precede bottoms)} for $x \in \{15, 16\}$ and
$1 \le j \le x/2$, the \emph{guard} $t_{x-2j} = 2M + 2j - x$ precedes
the \emph{bottom} $b_j = M + j$.
\end{enumerate}
$\OG(M)$ is \emph{infeasible} if no such order exists. (For $M \le 15$
some guard coincides with the bottom it protects --- $t_{x-2j} = b_j$
exactly when $M = x - j$, e.g.\ $t_{14} = b_1$ at $M = 15$ --- and the
gadget degenerates; we never use those scales. For $16 \le M \le 22$ a
guard can coincide with a \emph{different} bottom, which is harmless:
every attack pair is a genuine pair for all $M \ge 16$.)
\end{definition}

Three pieces of boundary arithmetic, each a one-line computation
(machine-checked exhaustively in the repository, script
\texttt{e96\_reduction\_check.py}), make the gadget exactly the trace of
the infinite problem on one block: for $x \in \{15,16\}$ the completion
$2b_j - x$ lies in $(M, 2M]$ precisely for $j \le x/2$, and then equals
the guard $t_{x-2j}$ (at $(x, j) = (16, 8)$ it is $2M$ itself, still
inside the half-open block) --- so (ii) lists \emph{all} in-block
completions of pairs $(x, b_j)$ and nothing else; all other completions
against $15$ and $16$ escape into the odd block above; and the
guard--bottom pair's own downward completion $2b_j - t_{x-2j} = x$ is the
attacking value itself, below the block, so the attack edges generate no
unmodeled in-block constraint.

\begin{theorem}[Order-gadget reduction; unconditional]\label{thm:ogred}
If $\OG(2^{2t-1})$ is infeasible for infinitely many $t \ge 4$, then
$\SA$ is not $3$-permutable.
\end{theorem}

\begin{proof}
Suppose $\pi$ is a $3$-AP-free permutation of $\SA$ and let
$P = \max(\pos(15), \pos(16))$; note $15, 16 \in B_4 \subseteq \SA$. Fix
$t \ge 4$ and put $M = 2^{2t-1}$, so that $(M, 2M] = B_{2t} \subseteq
\SA$ and $15, 16 < M$.

We claim that some bottom $b_j$ with $1 \le j \le 8$ has
$\pos(b_j) \le P$. Suppose not, and order the block by position:
$u \prec v :\!\iff \pos(u) < \pos(v)$. Constraint (i) is inherited from
$\pi$. For (ii), take $x \in \{15, 16\}$ and $1 \le j \le x/2$. The pair
$(x, b_j)$ is placed increasingly ($\pos(x) \le P < \pos(b_j)$, $x < M <
b_j$), and its completion $z = 2b_j - x = t_{x-2j}$ lies in
$(M, 2M] \subseteq \SA$ and differs from $b_j$ (else $M = x - j \le
15$). If $\pos(z) > \pos(b_j)$, then $(x, b_j, z)$ is a monotone $3$-AP
of $\pi$; hence $\pos(z) < \pos(b_j)$, i.e.\ $t_{x-2j} \prec b_j$. Thus
$\prec$ witnesses $\OG(M)$, contradicting infeasibility.

So for every $t \ge 4$ with $\OG(2^{2t-1})$ infeasible, some element of
$\{M+1, \dots, M+8\}$, $M = 2^{2t-1}$, occupies a position $\le P$.
These blocks are disjoint, so infinitely many such $t$ would place
infinitely many distinct values among the first $P$ positions ---
impossible.
\end{proof}

\begin{remark}[Chunk form; the crowns cannot defend every block]
\label{rem:ogchunk}
In the coordinates of Theorem~\ref{thm:chunk} the identical argument
reads: in any valid scheme $s$, for every $t \ge 4$ with
$\OG(2^{2t-1})$ infeasible, some bottom value in
$\{M+1, \dots, M+8\}$ of block $B_{2t}$ has stage
$\le \max(s(15), s(16))$. For if all eight had strictly larger stage,
the induced order on the block (by stage, then fiber position) would
satisfy $\OG(M)$: condition (i) holds for in-block triples directly by
(A) and (B), and each attack $t_{x-2j} \prec b_j$ is forced whenever
$s(x) < s(b_j)$ --- by (A) the guard's stage cannot strictly exceed the
bottom's, and at equal stages the forced-pair row
$s(x) < s(y) = s(z) \Rightarrow z \prec y$ of (B) applies. The overflow
step is then the finiteness of the fibers at stages
$\le \max(s(15), s(16))$; \emph{no} displacement normalization
(Lemma~\ref{lem:normal}) is needed anywhere in this route. This is the
precise sense in which the two \emph{crowns} $15$ and $16$ of
Section~\ref{sec:ladder} cannot defend every dyadic block: they would
have to be placed below the bottom-eight of all but finitely many
blocks. We use only the direction scheme $\Rightarrow$ order; the
converse (rebuilding a capped scheme from an $\OG$ witness, which would
make the crown-cap infeasibility at a horizon \emph{equivalent} to
$\OG$ infeasibility at its top scale) has been verified only at the
single scale $M = 128$, where the relevant family-MUS is single-block,
and is not used here.
\end{remark}

\begin{proposition}[Machine verification; Cadical195]\label{prop:ogmachine}
$\OG(M)$ is infeasible for every $M$ with $16 \le M \le 200$, and
for $M = 512$. In particular both tested scales of the dyadic family
$\{2^{2t-1}\}_{t \ge 4} = \{128, 512, 2048, \dots\}$ are infeasible.
The encodings assert (i) over all in-block triples and (ii) as unit
clauses; order axioms are handled by lazy transitivity refinement,
which is sound for infeasibility verdicts and was cross-validated
against full eager $O(n^3)$ transitivity encodings at $M = 40, 44$ and
--- as an end-to-end recheck --- at $M = 128$ (fresh eager instance,
$690{,}831$ clauses, UNSAT in about a second). The run at $M = 2048$
is unresolved (no verdict after $200+$ refinement rounds), and the
support growth reported below makes it unlikely that any finite sweep
can settle the general statement.
\end{proposition}

\subsection{Attack cores: locating the load-bearing axioms}

The hypothesis of Theorem~\ref{thm:ogred} demands $\OG$-infeasibility at
infinitely many dyadic scales. The following machine facts locate a
three-axiom core of the gadget that is already infeasible on exactly the
residue class containing the dyadic family; Section~\ref{sec:c3} proves
that core infeasibility by hand.

\begin{proposition}[Attack cores; machine-checked]\label{prop:cores}
For every swept $M$ (all of $40 \le M \le 100$, both parities, and
$M \in \{104, 108, 112, 120, 128, 150, 200\}$), constraint (i)
together with only the eleven attack units
$15\{1..7\} \cup 16\{1..4\}$ --- that is, $t_{15-2j} \prec b_j$ for
$j \le 7$ and $t_{16-2j} \prec b_j$ for $j \le 4$ --- is already
infeasible. On residue subclasses, sharper cores hold at every swept
scale: the four units $\{t_{13} \prec b_1,\ t_{11} \prec b_2,\
t_5 \prec b_5,\ t_{10} \prec b_3\}$ are infeasible with (i) exactly
when $M \equiv 0 \pmod 4$, and the three units
\[
\mathrm{C3} \;=\; \{\,t_5 \prec b_5,\quad t_3 \prec b_6,\quad
t_{10} \prec b_3\,\}
\]
are infeasible with (i) exactly when $M \equiv 0 \pmod 8$ (satisfiable
at every other swept residue, including $M \equiv 4 \bmod 8$).
\end{proposition}

Every scale of the dyadic family is $\equiv 0 \pmod 8$, so
Theorem~\ref{thm:ogred} needs only the infeasibility of
(i)${}\wedge{}\mathrm{C3}$ on the class $M \equiv 0 \pmod 8$. That is
exactly Theorem~\ref{thm:c3core} below.

\begin{remark}[Why a per-residue proof was to be expected]
The repository's parametric analysis (notes 27, 28, 30 and the
\texttt{e87}--\texttt{e113} experiment suite) had established three
structural facts that shaped the proof in Section~\ref{sec:c3}.
(1) Deletion-minimal certificate supports grow linearly with $M$
($\approx 8$ triples per unit of $M$), so no fixed finite list of
parametric AP-triples certifies infeasibility at all scales; a uniform
proof must use $\Theta(M)$-length mechanisms with $O(1)$ description ---
the ladder floods below are exactly that. (2) The forced-lemma layer
behind small-scale refutations is parity-locked (reverse-forced at
$M \not\equiv 0 \bmod 4$), so any uniform proof must stay within a
residue class; the dyadic family lives in $M \equiv 0 \pmod 8$.
(3) Fixed-denominator band-limit windows of the constraint system are
satisfiable at all scales, so the obstruction is genuinely asymptotic
and no compactness/limit argument over order types can carry it; the
proof must be a scale-uniform schema, verified rung by rung.
\end{remark}

\section{The C3 core theorem and the main theorem}\label{sec:c3}

\begin{theorem}[Main theorem]\label{thm:main}
$\SA = \bigcup_{k \text{ even}} (2^{k-1}, 2^k]$ is not $3$-permutable:
it admits no permutation (of order type $\omega$) free of monotone
$3$-term arithmetic progressions. In particular the canonical dyadic
partition $\Z^+ = \SA \sqcup S_B$ is not a solution of the
Erd\H{o}s--Graham two-set problem: any partition witnessing a YES answer
to problem \#197, if one exists, must differ from the dyadic one.
\end{theorem}

\begin{proof}
Every scale of the dyadic family satisfies $M = 2^{2t-1} \equiv 0
\pmod 8$ and $M \ge 128$ for $t \ge 4$. By Theorem~\ref{thm:c3core}
below, the C3 core --- hence a fortiori the full gadget $\OG(M)$, whose
attack list (Definition~\ref{def:og}(ii)) contains the three axioms of
$\mathrm{C3}$ as the instances $(x, j) = (15, 5), (15, 6), (16, 3)$,
and whose constraint (i) is AP-freeness
--- is infeasible at every such scale. Theorem~\ref{thm:ogred} applies.
\end{proof}

The remainder of this section proves the core statement:

\begin{theorem}[C3 core]\label{thm:c3core}
For every $M \equiv 0 \pmod 8$ with $M \ge 16$, no AP-free linear order
of the interval $(M, 2M]$ satisfies the three axioms
\[
\mathrm{C3} = \{\, A_1\colon t_5 \prec b_5,\quad
A_2\colon t_3 \prec b_6,\quad A_3\colon t_{10} \prec b_3 \,\}.
\]
\end{theorem}

Here \emph{AP-free} means condition (i) of Definition~\ref{def:og}: no
arithmetic progression $a<b<c$ inside the block occurs in monotone
position order. The proof is a small toolkit of ladder lemmas
(\S\ref{sec:toolkit}) followed by two forcing theorems
(\S\ref{sec:l1}--\S\ref{sec:flip}). Every lemma application in the
proofs has been machine-verified step by step at every
$M \equiv 0 \pmod 4$ (resp.\ $\equiv 0 \pmod 8$) in $[12, 400]$ plus
$M = 512, 1024$, and independently cross-validated; see
Remark~\ref{rem:c3machine}.

\subsection{The ladder toolkit}\label{sec:toolkit}

AP-freeness is the \emph{midpoint-extremal rule}: on every in-block AP
$(a, b, c)$, $c = 2b - a$, the midpoint $b$ either precedes both
endpoints or follows both. We use it as four unit rules (plus
transitivity throughout):
\[
\text{R1: } a \prec b \Rightarrow c \prec b, \qquad
\text{R3: } c \prec b \Rightarrow a \prec b, \qquad
\text{R2: } b \prec c \Rightarrow b \prec a, \qquad
\text{R4: } b \prec a \Rightarrow b \prec c.
\]
Write $m_0 = 3M/2$ (the arithmetic midpoint of the block: $b_j + t_j =
3M = 2m_0$ for all $j$, so every pair $(b_j, t_j)$ mirrors through
$m_0$). A \emph{$d$-ladder} is a maximal run $w_0, w_1, \dots$ of
values of $(M, 2M]$ in arithmetic progression with difference $d$; the
$d = 2$ ladders are the two parity classes, and the $d = 4$ ladders on
odd values are \emph{class A} $= \{v \equiv M + 1 \bmod 4\}$ and
\emph{class B} $= \{v \equiv M + 3 \bmod 4\}$.

\begin{lemma}[Zigzag]\label{lem:zigzag}
Let $w_0, \dots, w_r$ be consecutive rungs of a $d$-ladder and suppose
$w_e \prec w_{e'}$ for some $|e - e'| = 1$. Then every rung $w_i$ with
$i \equiv e \pmod 2$ leads both its neighbors ($w_i \prec w_{i \pm 1}$).
\end{lemma}

\begin{proof}
Any three consecutive rungs form an AP with the middle rung as
midpoint, so each rung either leads or trails both neighbors. The seed
makes $w_e$ a leader. If $w_i$ leads, then from $w_i \prec w_{i+1}$,
R1 on $(w_i, w_{i+1}, w_{i+2})$ gives $w_{i+2} \prec w_{i+1}$, and R4
on $(w_{i+1}, w_{i+2}, w_{i+3})$ gives $w_{i+2} \prec w_{i+3}$: $w_{i+2}$
leads. Downward is the mirror argument.
\end{proof}

\begin{lemma}[Phase dichotomy]\label{lem:phase}
In any linear order, every $d$-ladder is globally in one of its two
\emph{zigzag phases}: either all even-index rungs lead both neighbors,
or all odd-index rungs do. Consequently the \emph{leader set} of the
ladder is one of its two half-classes modulo $2d$, adjacent rungs
strictly alternate leader/trailer, and a proof may case-split on the
phase and detach any conclusion derived in both branches.
\end{lemma}

\begin{proof}
The order orients the adjacent pair $(w_0, w_1)$ one way or the other;
apply Lemma~\ref{lem:zigzag} to that seed.
\end{proof}

\begin{lemma}[Transfer lock]\label{lem:transfer}
Let $M$ be even. On the odd $d=2$ ladder $w_i = M + 1 + 2i$
($0 \le i \le M/2 - 1$) the four odd C3 values sit at $b_3 = w_1$,
$b_5 = w_2$, $t_5 = w_{M/2-3}$, $t_3 = w_{M/2-2}$, and
Lemma~\ref{lem:zigzag} locks the pair orientations together:
\[
M \equiv 0 \ (\mathrm{mod}\ 4):\quad b_5 \prec b_3
\iff t_3 \prec t_5;
\qquad
M \equiv 2 \ (\mathrm{mod}\ 4):\quad b_5 \prec b_3 \iff t_5 \prec t_3.
\]
\end{lemma}

\begin{proof}
Each orientation of an adjacent pair seeds Lemma~\ref{lem:zigzag}; read
off whether the other pair's left member is a leader. At
$M \equiv 0 \pmod 4$: $b_5 \prec b_3$ (i.e.\ $w_2 \prec w_1$) makes
even indices leaders, and $M/2 - 2$ is even, so $t_3 = w_{M/2-2}$
leads its neighbor $t_5$. Conversely $t_3 \prec t_5$ seeds even-index
leaders and $w_2 = b_5$ leads $w_1 = b_3$. The reverse orientations
force the reverse conclusions, giving the biconditional; at
$M \equiv 2 \pmod 4$ the parity of $M/2$ shifts the lock by one rung.
\end{proof}

\begin{lemma}[Flood]\label{lem:flood}
Fix $g \in \{2, 4\}$ and a $g$-class $C$ of $(M, 2M]$ (a parity class
for $g = 2$; one of the odd mod-$4$ classes for $g = 4$) whose
$d = g$ ladder is in a definite zigzag phase with leader set $L$. Let
$c \in (M, 2M]$ be a \emph{center} with $c \equiv C + g/2 \pmod g$, so
that the mirror pairs $(c - e,\, c,\, c + e)$ with $e \equiv g/2
\pmod g$ have both members in $C$ and form APs with midpoint $c$; call
$e$ \emph{admissible} when both $c \pm e \in (M, 2M]$. Suppose some
admissible $e_0$ carries a seed relation between $c$ and a member
$v \in \{c \pm e_0\}$. Then:
\begin{itemize}
\item (outward) if $c \prec v$, then $c \prec w$ for \emph{every}
$w \in C$ with $2c - w \in (M, 2M]$;
\item (inward) if $v \prec c$, then $w \prec c$ for every such $w$.
\end{itemize}
Moreover the conclusion holds in both zigzag phases of the $C$-ladder
(\emph{phase-blindness}), so by Lemma~\ref{lem:phase} it may be
detached after a two-branch case split whenever the seed is available
in both branches.
\end{lemma}

\begin{proof}
First, at each admissible $e$ exactly one of $c + e$, $c - e$ is a
leader: they differ by $2e \equiv g \pmod{2g}$, so they lie in the two
different half-classes mod $2g$ of $C$, and $L$ is one of them ---
whichever phase holds. Second, in a zigzag, adjacent rungs alternate
and each leader leads both its neighbors.

\emph{Seed.} From the relation at $v$, the mirror rule on the AP
$(c - e_0, c, c + e_0)$ (R2/R4 outward, R1/R3 inward) gives the
same-direction relation at the other member: both members of pair
$e_0$ are related to $c$ in the flood direction.

\emph{Outward, } $e \to e + g$: let $x \in \{c \pm e\}$ be the leader
of pair $e$. Its outward ladder-neighbor $x'$ (same side,
$|x' - c| = e + g$) satisfies $x \prec x'$; with $c \prec x$ known,
transitivity gives $c \prec x'$, and the mirror rule floods
$2c - x'$.

\emph{Outward, } $e \to e - g$: let $w \in \{c \pm (e - g)\}$ be the
trailer of pair $e - g$. Its outward neighbor $w'$ (distance $e$, same
side) is a leader and leads it: $w' \prec w$; with $c \prec w'$ known,
$c \prec w$, and the mirror rule floods the other member.

\emph{Inward, } $e \to e + g$: let $x$ be the leader of pair $e + g$;
it leads its inward neighbor $x_{\mathrm{in}}$ (distance $e$, same
side): $x \prec x_{\mathrm{in}} \prec c$, and the mirror rule floods
the other member. \emph{Inward, } $e \to e - g$: let $x$ be the leader
of pair $e - g$; it leads its outward neighbor (distance $e$):
$x \prec x_{\mathrm{out}} \prec c$; mirror floods the other member.

The admissible $e$ form an interval of the residue class $g/2 \bmod g$,
so the two inductions from $e_0$ reach every admissible $e$. Every step
used only R1--R4 on genuine APs, transitivity, and zigzag edges of the
given phase; the choice of side at each step is forced by the phase but
\emph{exists} in both phases.
\end{proof}

Three instances of Lemma~\ref{lem:flood} are used below.
\textbf{POLAR} ($g = 2$, center $m_0$,
class the odds; $M \equiv 0 \bmod 4$ makes $m_0$ even): seeded by the
comparison of $m_0$ with $t_5$ (mirror $b_5$), it floods $m_0$ against
\emph{all} odd values (the mirror of $b_l$ is $t_l$). The odd ladder's
phase is pinned by the ambient hypotheses, so no dichotomy is needed.
\textbf{P2-floods} ($g = 2$, an odd center $c$ near $m_0$, class the
evens): seeded by $c$ vs.\ $m_0$ at $e_0 = |c - m_0|$ (supplied by
POLAR); the even ladder's phase is unknown --- phase dichotomy, both
branches. \textbf{G4-floods} ($g = 4$, center $c \equiv \text{class} +
2 \pmod 4$, class A or B): seeded at $e = 2$ by the odd ladder's zigzag
edges between $c$ and $c \pm 2$; the $d = 4$ ladder's phase is unknown
--- phase dichotomy, both branches.

The \emph{mod-8 lock} of the whole problem sits in the center condition
of the G4-floods: at $M \equiv 0 \pmod 8$ the two odd neighbors of
$m_0$ satisfy $m_0 - 1 \equiv 3$ and $m_0 + 1 \equiv 1 \pmod 4$, so
$m_0 - 1$ is a G4-center for class A and $m_0 + 1$ for class B. At
$M \equiv 4 \pmod 8$ both congruences fail (and the odd-ladder leader
statuses at $m_0 \pm 1$ invert), the flood seeds do not exist, and the
proofs below evaporate --- matching the machine fact that
(i)${}\wedge{}\mathrm{C3}$ is \emph{satisfiable} there
(Proposition~\ref{prop:cores}).

\subsection{Layer 1: two axioms force the odd-pair orientations}
\label{sec:l1}

\begin{theorem}[L1]\label{thm:l1}
Let $M \equiv 0 \pmod 4$, $M \ge 12$, and let $\prec$ be an AP-free
order of $(M, 2M]$ satisfying $A_2\colon t_3 \prec b_6$ and
$A_3\colon t_{10} \prec b_3$. Then $b_5 \prec b_3$ and $t_3 \prec t_5$.
\end{theorem}

\begin{proof}
It suffices to refute $S\colon b_3 \prec b_5$; given $b_5 \prec b_3$,
Lemma~\ref{lem:transfer} supplies $t_3 \prec t_5$. Assume $A_2, A_3,
S$. By Lemma~\ref{lem:zigzag} (seed $S$) the odd ladder's leaders are
the offsets $\equiv 3 \pmod 4$. Define
\[
c^{*} =
\begin{cases} m_0 + 1, & M \equiv 0 \ (\mathrm{mod}\ 8)\\
m_0 - 1, & M \equiv 4 \ (\mathrm{mod}\ 8)\end{cases}
\qquad
c^{**} =
\begin{cases} m_0 - 1, & M \equiv 0 \ (\mathrm{mod}\ 8)\\
m_0 + 1, & M \equiv 4 \ (\mathrm{mod}\ 8)\end{cases}
\]
so that $c^{*} \equiv 1 \pmod 4$ is a G4-center for class B whose odd
neighbors $c^{*} \pm 2$ ($\equiv 3 \bmod 4$) are odd-ladder leaders,
and $c^{**} \equiv 3 \pmod 4$ is a G4-center for class A and is itself
an odd-ladder leader. Split on the comparison $(m_0, t_5)$.

\emph{Case I: $t_5 \prec m_0$.} POLAR-inward: every odd value
$\prec m_0$. Then:
(1) $b_3 \prec c^{*}$ by the G4-inward flood at $c^{*}$ over class B
--- seeds $c^{*} \pm 2 \prec c^{*}$ (the leaders $c^{*} \pm 2$ lead
their trailing neighbor $c^{*}$); $b_3 \in$ class B with mirror
$2c^{*} - b_3 \in \{t_1, t_5\}$ in the block; both $d{=}4$ phases.
(2) $c^{*} \prec t_{10}$ by the P2-outward flood at $c^{*}$ over the
evens --- seed $c^{*} \prec m_0$ (POLAR, $e_0 = 1$); $t_{10}$ even
with mirror $2c^{*} - t_{10} \in \{M{+}8, M{+}12\}$ in the block; both
even phases.
(3) $A_3$: $t_{10} \prec b_3$. Together: the $3$-cycle
$b_3 \prec c^{*} \prec t_{10} \prec b_3$, a contradiction.

\emph{Case II: $m_0 \prec t_5$.} POLAR-outward: $m_0 \prec$ every odd
value. Then:
(1) $c^{**} \prec t_3$ by the G4-outward flood at $c^{**}$ over class
A --- seeds $c^{**} \prec c^{**} \pm 2$ ($c^{**}$ is an odd-ladder
leader); $t_3 \in$ class A with mirror $2c^{**} - t_3 \in
\{b_1, b_5\}$; both phases.
(2) $b_6 \prec c^{**}$ by the P2-inward flood at $c^{**}$ over the
evens --- seed $m_0 \prec c^{**}$ (POLAR, $e_0 = 1$); $b_6$ even with
mirror $2c^{**} - b_6 \in \{2M{-}8, 2M{-}4\}$; both phases.
(3) $A_2$: $t_3 \prec b_6$. Together: the $3$-cycle
$c^{**} \prec t_3 \prec b_6 \prec c^{**}$, a contradiction.

Both cases are contradictory, so $S$ is impossible.
\end{proof}

Note that Case I consumes only $A_3$ and Case II only $A_2$; the
boundary arithmetic (mirrors in the block, seed admissibility) holds
down to $M = 12$.

\subsection{The flip: the third axiom is anti-forced}\label{sec:flip}

\begin{theorem}[FLIP]\label{thm:flip}
Let $M \equiv 0 \pmod 8$, $M \ge 16$, and let $\prec$ be an AP-free
order of $(M, 2M]$ satisfying $A_2$, $A_3$ and $b_5 \prec b_3$. Then
$t_5 \prec b_5$ is impossible: $b_5 \prec t_5$ ($= \lnot A_1$) is
forced.
\end{theorem}

\begin{proof}
Assume $A_1\colon t_5 \prec b_5$ for contradiction. By
Lemma~\ref{lem:zigzag} (seed $b_5 \prec b_3$) the odd ladder's leaders
are the offsets $\equiv 1 \pmod 4$. Since $M \equiv 0 \pmod 8$:
$m_0 - 1 \equiv 3 \pmod 4$ is a G4-center for class A whose odd
neighbors $m_0 + 1$, $m_0 - 3$ (offsets $\equiv 1 \bmod 4$) are
odd-ladder leaders; $m_0 + 1 \equiv 1 \pmod 4$ is a G4-center for
class B and is itself an odd-ladder leader. Split on $(m_0, t_5)$.

\emph{Case I: $t_5 \prec m_0$.} POLAR-inward: every odd ${}\prec m_0$.
Then:
(1) $m_0 - 1 \prec t_{10}$ [P2-outward flood at $m_0 - 1$ over the
evens; seed $m_0 - 1 \prec m_0$ (POLAR, $e_0 = 1$); mirror of
$t_{10}$ is $M + 8$; both phases].
(2) $m_0 - 1 \prec b_3$ [(1) $+ A_3$].
(3) $m_0 - 1 \prec t_5$ [R4 on the AP $(b_3,\, m_0 - 1,\, t_5)$:
$b_3 + t_5 = 3M - 2 = 2(m_0 - 1)$].
(4) $m_0 - 1 \prec b_5$ [(3) $+ A_1$].
(5) $b_5 \prec m_0 - 1$ [G4-inward flood at $m_0 - 1$ over class A;
seeds $m_0 + 1 \prec m_0 - 1$ and $m_0 - 3 \prec m_0 - 1$ (leader
edges); $b_5 \in$ class A at pair distance $M/2 - 6 \equiv 2 \pmod 4$
with mirror $t_7$ in the block; both phases].
(4) and (5) form a $2$-cycle, a contradiction.

\emph{Case II: $m_0 \prec t_5$.} POLAR-outward: $m_0 \prec$ every odd.
Then:
(1) $b_6 \prec m_0 + 1$ [P2-inward flood at $m_0 + 1$ over the evens;
seed $m_0 \prec m_0 + 1$ (POLAR, $e_0 = 1$); mirror of $b_6$ is
$2M - 4$; both phases].
(2) $t_3 \prec m_0 + 1$ [$A_2$ $+$ (1)].
(3) $b_5 \prec m_0 + 1$ [R3 on the AP $(b_5,\, m_0 + 1,\, t_3)$:
$b_5 + t_3 = 3M + 2 = 2(m_0 + 1)$].
(4) $t_5 \prec m_0 + 1$ [$A_1$ $+$ (3)].
(5) $m_0 + 1 \prec t_5$ [G4-outward flood at $m_0 + 1$ over class B;
seeds $m_0 + 1 \prec m_0 + 3$ and $m_0 + 1 \prec m_0 - 1$ ($m_0 + 1$
is an odd-ladder leader); $t_5 \in$ class B at pair distance
$M/2 - 6 \equiv 2 \pmod 4$ with mirror $b_7$; both phases].
(4) and (5) form a $2$-cycle, a contradiction.
\end{proof}

Again Case I consumes $A_3$ and Case II consumes $A_2$, while $A_1$ is
consumed symmetrically in both (step 4). The proof visibly fails at
$M \equiv 4 \pmod 8$ in both cases: the centers $m_0 \mp 1$ land in
the wrong mod-$4$ classes for their G4-floods and the odd-ladder
leader statuses at $m_0 \pm 1$ invert, so neither seed set exists.

\begin{proof}[Proof of Theorem~\ref{thm:c3core}]
$A_2 + A_3$ force $b_5 \prec b_3$ (Theorem~\ref{thm:l1};
$M \equiv 0 \bmod 8 \subseteq 0 \bmod 4$). Then $A_1$ contradicts
Theorem~\ref{thm:flip}.
\end{proof}

\begin{remark}[Machine verification of the schemas]\label{rem:c3machine}
Although the proofs above are complete hand proofs, each is a
\emph{schema} whose instances at a given scale involve $\Theta(M)$
ladder steps; all have been machine-executed with strict per-step
assertions (every AP's membership and arithmetic, every rule pattern,
every leader/trailer and residue claim, every mirror bound, both
branches of every phase dichotomy): Theorem~\ref{thm:l1} at every
$M \equiv 0 \pmod 4$ in $[12, 400]$ plus $512, 1024$ ($100$ scales),
Theorem~\ref{thm:flip} at every $M \equiv 0 \pmod 8$ in $[16, 400]$
plus $512, 1024$ ($51$ scales), and the $M \equiv 4 \pmod 8$
inapplicability at $35$ scales
(script \texttt{e113\_c3\_hand\_proof.py}). An independent closure
engine (plain R1--R4 $+$ transitivity fixpoint, no knowledge of the
schemas) refutes every branch of both case trees at $51$ resp.\ $27$
scales up to $512$ (\texttt{e113b\_closure\_crossval.py}); and direct
SAT checks of the two theorem \emph{statements} at fresh scales never
touched by the discovery loop ($M = 148, 212, 264, 328$ UNSAT;
$M = 268, 332$, $\equiv 4 \bmod 8$, SAT) concur
(\texttt{e114\_theorem\_spotcheck.py}). The earlier exhaustive base
ledgers (AP${}+{}$C3 UNSAT at every $M \equiv 0 \bmod 8$ in
$[16, 256]$ and at $512$) are now corroboration, not a dependency.
\end{remark}

\begin{remark}[What remains open]\label{rem:c3open}
The main theorem needs nothing beyond the class $M \equiv 0 \pmod 8$.
The full order gadget is nevertheless conjectured infeasible at
\emph{every} $M \ge 16$ (machine-verified for $16 \le M \le 200$ and
$M = 512$, Proposition~\ref{prop:ogmachine}); at residues
$\not\equiv 0 \bmod 8$ the C3 core is satisfiable and other attack
subsets take over (Proposition~\ref{prop:cores}), so a proof would
need per-residue analogues of Theorems~\ref{thm:l1}--\ref{thm:flip}
built from the same flood toolkit. This statement is not needed for
Theorem~\ref{thm:main}. The Erd\H{o}s--Graham problem \#197 itself
remains open: Theorem~\ref{thm:main} eliminates the canonical
candidate partition (and with it the natural route to a YES), but
does not decide whether some other two-set partition works.
\end{remark}

\section{Historical note: suffix-stacked deferral}

The deferral-state experiments (states exist at $M=32,64,128$; class
coherence; sink algebra) are subsumed by the chunk coordinates: prefix-nested
semantics permit top-half deferral that rigid chunks forbid, which is why the
law's witnesses evade Lemma~\ref{lem:bottomhalf} while no closed-form law
survived depth $4$.

\section{Remarks and open problems}

% TODO: the LV density program; measurement data (g-curves); the exact
% characterization discussion.

\end{document}